% =====================================================================
% DPRNet -- Springer SUBMISSION build for The Visual Computer (TVCJ).
% Created 2026-08-01.
% Updated for TVCJ first submission on 2026-09-29.
%
% WHY a standard article class instead of sn-jnl.cls:
%   The Springer Nature class file (sn-jnl.cls) is distributed only via
%   Springer's own site / Overleaf and is NOT on CTAN; it could not be
%   fetched in this environment. Per the official author guidelines
%   (verified 2026-08-01):
%     - The Visual Computer accepts editable LaTeX source and has no
%       journal-specific official template at first submission.
%     - Complete editable source files are required at every submission.
%   A single-column 12pt article source therefore satisfies "editable
%   source" for first submission. If the paper reaches revision/acceptance,
%   re-typeset into sn-jnl.cls ([iicol]) then.
%
% Reuses the SAME shared sources as the CVPR (main.tex), MDPI
% (main_mdpi.tex), and Elsevier (main_els.tex) builds:
%   sections/*  tables/*  figures/*  references.bib
%
% Compile:
%   pdflatex main_springer ; bibtex main_springer ; pdflatex main_springer ; pdflatex main_springer
% =====================================================================
\documentclass[12pt]{article}

% ---- fonts: Latin Modern provides Type1 CM + TC (TS1) glyphs, so pdflatex does
% NOT need to generate pk bitmap fonts (tcrm*/ecrm*). Required here because the
% default article Computer Modern triggers on-the-fly pk generation, which this
% sandbox blocks; the Times-based els/mdpi/cvpr builds never hit it.
\usepackage{lmodern}
\usepackage[T1]{fontenc}

% ---- page geometry: single column, generous margins (review-friendly) ----
\usepackage[a4paper,margin=1in]{geometry}
\usepackage{setspace}
\onehalfspacing
\setlength{\emergencystretch}{3em}
\usepackage[numbers,sort&compress]{natbib}   % Springer Nature Basic style is numeric
\bibliographystyle{unsrtnat}                  % numeric, cited-order (Springer-compatible;
                                              % sn-jnl.cls not on CTAN in this env)

\input{preamble_els}   % shared extra packages for figures/tables (amsmath, booktabs, multirow, float, pifont, pgfplots, tikz, dprblue/dpramber, \TODO ...)
\hypersetup{
  pdfauthor={Hainan Huang and Xiang Li},
  pdftitle={Dynamic Prototype Routing for Lightweight Image Super-Resolution}
}

\title{\bfseries Dynamic Prototype Routing for Lightweight Image Super-Resolution}

\author{
  Hainan Huang\textsuperscript{*} and Xiang Li\\[4pt]
  \normalsize School of Computer Science\\
  \normalsize China University of Geosciences (Wuhan)\\
  \normalsize Future City Campus, No.\ 68 Jincheng Street\\
  \normalsize East Lake High-tech Development Zone, Wuhan, Hubei 430078, China\\[3pt]
  \small \textsuperscript{*}Corresponding author:
  \texttt{huanghainan.7@cug.edu.cn}; Xiang Li: \texttt{lixiang@cug.edu.cn}
}

\date{}

\begin{document}
\maketitle

\begin{abstract}
\noindent
Lightweight single image super-resolution (SR) must balance reconstruction quality with
model cost on resource-constrained platforms. Content-aware routing improves this balance
by grouping semantically related tokens, yet the widely used center-based router---the
Token Aggregation Block (TAB) of CATANet---relies on cross-batch exponential-moving-average
centers, hard assignment labels, and a coupled prototype update/confirmation step, tying
routing to a drifting historical state and discarding assignment confidence. We revisit
token routing itself and propose \emph{Dynamic Prototype Routing} (DPR), a confidence-aware,
drop-in replacement for the TAB routing that keeps the CATANet backbone and tensor interface
unchanged. DPR builds per-image prototypes directly from the current tokens, decouples
prototype generation from confirmation through a learnable query bank, matches tokens with a
bounded learnable temperature, and carries the resulting assignment confidence into token
ordering, global gating, and a soft fallback, with a usage-balancing regularizer that
discourages prototype-slot collapse. Trained on DIV2K and evaluated on five standard
benchmarks at $\times2/\times3/\times4$, DPRNet matches CATANet-level accuracy with small
gains on most benchmark cells while staying in the lightweight regime; at $\times4$, for
example, it obtains $26.89$\,dB on Urban100 with $52.14$\,G Multi-Adds, fewer than the
lightweight Transformer baselines SwinIR-light and SRFormer-light under the same input
convention. Three-seed ablations isolate each design choice: the dynamic prototypes
outperform a history-center EMA variant under an identical protocol, and the balancing loss
makes prototype usage more balanced and more stable across seeds. DPR thus offers a reusable,
interpretable alternative to center-based token routing rather than a larger backbone.
\end{abstract}

\noindent\textbf{Keywords:} image super-resolution \textperiodcentered\ lightweight
super-resolution \textperiodcentered\ content-aware routing \textperiodcentered\ dynamic
prototype routing \textperiodcentered\ vision transformer \textperiodcentered\ self-attention

\bigskip

% =====================================================================
% Section: Introduction (with Related Work)  (outline.md §1)
% Draft v1 (2026-06-10). Related Work is integrated into the Introduction
% following the CS/engineering SCI body convention (outline.md).
% Citation keys refer to paper/references.bib.
% This draft contains NO quantitative result assertions; all comparative
% statements are qualitative (motivation / positioning) and are backed by
% the experiments referenced in plan/review/method-experiment-traceability.md.
% =====================================================================
\section{Introduction}
\label{sec:intro}

Single image super-resolution (SR) aims to reconstruct a high-resolution (HR)
image from a low-resolution (LR) observation, a problem that underpins applications
from mobile photography to display rendering. Before deep learning, SR was often
formulated with hand-crafted image priors, such as internal self-similarity
across scales~\cite{glasner2009single}. Since the seminal deep
model SRCNN~\cite{dong2014srcnn}, accuracy has improved through deeper,
residual, and feature-reuse networks, including
VDSR~\cite{kim2016vdsr}, RDN~\cite{zhang2018rdn}, RCAN~\cite{zhang2018rcan},
and EDSR~\cite{lim2017edsr}. A recent survey summarizes this progression across
architectures, losses, datasets, and metrics~\cite{moser2023survey}. The
corresponding growth in parameters and computation, however, makes many
high-capacity models impractical for mobile and edge devices, where memory,
latency, and energy are tightly constrained. This tension has motivated a sustained
line of \emph{lightweight} SR, which seeks a favorable
trade-off between reconstruction quality and model cost.

\paragraph{Lightweight CNN-based SR.}
Efficient CNN designs reduce cost by moving computation to low-resolution
features, reusing parameters, simplifying feature distillation, or replacing
heavy operators with compact mixing. CARN employs cascading
residual connections with group convolutions~\cite{ahn2018carn}; IMDN introduces
information multi-distillation~\cite{hui2019imdn}; RFDN reformulates this idea
with residual feature distillation blocks~\cite{liu2020rfdn}; and RLFN further
simplifies local feature extraction for efficient SR~\cite{kong2022rlfn}. More
recent lightweight designs also explore shuffle-based convolutional
mixing~\cite{sun2022shufflemixer} and dual-regression
constraints~\cite{guo2022dualregression}.
Efficient-SR challenges reflect the same deployment pressure by ranking models
with accuracy, runtime, parameters, and FLOPs considered together~\cite{ren2024ntire}.
For the present work, the main limitation of these methods is not efficiency but
context: locally connected convolutional operators cannot directly compare spatially
distant yet semantically similar regions without stacking many layers.

\paragraph{Lightweight Transformer-based SR.}
Long-range modeling tackles the context limitation from a different direction.
Non-local operators aggregate features across distant positions~\cite{wang2018nonlocal},
and Transformer-based restoration extends this idea with learned attention.
Swin Transformer provides the shifted-window
operator that later restoration models adopt for tractable attention~\cite{liu2021swin}.
SwinIR introduces window-based attention for image restoration~\cite{liang2021swinir};
ELAN streamlines long-range attention to cut redundant computation~\cite{zhang2022elan};
SRFormer revisits the window--channel trade-off with permuted
self-attention~\cite{zhou2023srformer}; and HAT activates more pixels to enlarge the
effective receptive field~\cite{chen2023hat}. Efficient variants further mix local
operators with omni-dimensional attention~\cite{wang2023omnisr} or n-gram
context~\cite{choi2023ngswin}, adopt mixed Transformer
blocks~\cite{zheng2023emt}, or integrate features hierarchically from coarse to
fine~\cite{liu2022hpinet}. More recent lightweight
Transformers push this efficiency further: GRFormer introduces grouped residual
self-attention to trim redundant parameters and computation~\cite{li2024grformer},
LMLT applies multi-level attention with per-head feature sizes to widen the
effective context~\cite{kim2025lmlt}, and Compacter shares a single projection
across spatial and channel attention for compact restoration~\cite{wu2024compacter}.
ATD takes a step toward content adaptivity by augmenting local windows with an
adaptive token dictionary~\cite{zhang2024atd}. Yet all of these still bind token
interaction to a predefined spatial partition: semantically similar but spatially
distant tokens interact only when the window grows or an extra global path is
added, which motivates organizing tokens by content instead of position.

\paragraph{Content-adaptive routing and token clustering.}
A complementary direction routes computation or interaction by \emph{content}
rather than by a fixed spatial layout. One thread allocates compute where
reconstruction is hard: PCSR classifies each pixel and dispatches it to an
upsampler of matching capacity~\cite{jeong2024pcsr}, while SalDRN treats visual
saliency as a proxy for region-level difficulty and selects feature paths of
appropriate depth accordingly~\cite{wu2023saldrn}. Pushing this idea to
real-world SR, MoR-DASR casts restoration as a mixture of experts whose
degradation-aware gate activates only relevant ranks and relies on a
load-balancing loss to stop expert usage from collapsing onto a few
branches~\cite{he2026mordasr}. A second thread instead groups \emph{tokens} by
content, so that distant but semantically related regions interact within a
compact attention---an idea rooted in non-local and self-exemplar SR that mines
recurring patches within an image~\cite{mei2020csnln} or transfers texture from a
reference~\cite{yang2020ttn}. SPIN clusters tokens into superpixels for
lightweight interaction~\cite{zhang2024spin}, and CATANet, the most direct
predecessor of our work, introduces a Token Aggregation Block (TAB) that keeps a
small set of routing centers, assigns each token to its nearest center, reorders
tokens by cluster label, and aggregates the reordered sequence with an
intra/inter-group attention operator~\cite{liu2025catanet}. This content routing
attains strong accuracy at a low budget; crucially, the load-balancing concern
raised for expert routing reappears here as a risk of \emph{prototype collapse}
that any learnable token router must also control.

\paragraph{Limitations of center-based routing.}
Despite its effectiveness, the TAB routing mechanism has four limitations.
(i) Its routing centers are maintained across batches through a $k$-means--style
iteration with an exponential moving average (EMA) buffer, so the centers are
driven by a cross-batch historical state rather than strictly by the current
image and adapt only slowly to per-image semantics. (ii) The token-to-center
relation is a hard \texttt{argmax} label that discards the \emph{confidence} of
the assignment, treating a token that clearly belongs to one cluster identically
to an ambiguous one on a boundary. (iii) The reordering clusters tokens only by
label, with no intra-cluster purity constraint, so high- and low-confidence tokens
are interleaved within a group during aggregation. (iv) Prototype \emph{generation}
and \emph{confirmation} are entangled in a single iterative update, which weakens
interpretability and stability. These observations suggest that a routing module
built directly from the current input, that preserves and exploits assignment
confidence, and that separates prototype generation from confirmation could
improve content routing without sacrificing the lightweight regime---provided it
also guards against the prototype collapse that a learnable router can induce, in
the same spirit as the load balancing used for expert routing~\cite{he2026mordasr}.

\paragraph{Our approach.}
We propose \emph{Dynamic Prototype Routing} (DPR), a drop-in replacement for the
TAB routing that keeps the CATANet backbone and its external tensor interface
unchanged while remaining in a comparable lightweight operating regime. Rather than
enlarging the network, our contribution is at the level of the routing mechanism:
DPR turns center-based token routing into a confidence-aware process. It builds
semantic prototypes directly from the current tokens via a soft assignment, confirms
them with a learnable query bank that separates prototype generation from
confirmation, matches tokens back to the confirmed prototypes with a learnable
temperature, and propagates an explicit per-token confidence through ordering,
gating, and a soft fallback, with a usage-balancing regularizer that guards against
prototype collapse. The resulting network, \textbf{DPRNet}, is trained on DIV2K and
evaluated on five standard benchmarks at $\times2/\times3/\times4$, where it maintains
CATANet-level reconstruction quality with small gains on most benchmark cells and a
comparable lightweight cost profile. Just as importantly, each mechanism is isolated
by a controlled, multi-seed ablation and the learned routing is made interpretable
through per-token confidence and prototype-usage diagnostics, so the improvement is
attributable to the routing design itself rather than to added capacity.

\paragraph{Contributions.}
The main contributions are as follows.
\begin{itemize}
  \item \textbf{C1 -- Dynamic semantic prototypes.} We generate routing prototypes
  from the current tokens by a soft, assignment-weighted aggregation, removing the
  cross-batch EMA center buffer so that prototypes follow per-image semantics
  (Section~\ref{sec:method_dpr}, Eqs.~\eqref{eq:assign}--\eqref{eq:proto_content}).
  \item \textbf{C2 -- Decoupled prototype confirmation.} We separate prototype
  \emph{generation} from \emph{confirmation} by refining the content prototypes
  with a learnable query bank through a single gated cross-attention, improving the
  stability and interpretability of the routing slots
  (Eq.~\eqref{eq:proto_refine}).
  \item \textbf{C3 -- Confidence-aware routing.} We explicitly retain the token
  assignment confidence and use it in three places: a composite ordering key that
  sorts high-confidence tokens first within each cluster, a confidence-modulated
  global gate, and a soft prototype fallback for uncertain tokens
  (Eqs.~\eqref{eq:sortkey}, \eqref{eq:iasa_gate}, \eqref{eq:soft_fallback}).
  \item \textbf{C4 -- Stabilized routing.} We introduce a learnable routing
  temperature together with a prototype usage-balancing regularizer to mitigate
  collapse onto a few prototypes and keep slot usage spread out
  (Eqs.~\eqref{eq:scores}, \eqref{eq:balance}).
\end{itemize}
Each design element is exposed as an independent switch that defaults to the full
model. Turning these switches off gives the corresponding ablated variants, which
allows the contributions to be isolated cleanly in the ablation study
(Section~\ref{sec:ablation}).

\input{sections/2_method}
\input{figures/fig1_architecture}
\input{figures/fig2_dpr}

% ---- Experiments ----
\input{sections/3_experiments_setup}
\input{sections/3_experiments_comparison}
\input{tables/table1_main_comparison}
\input{tables/table2_efficiency}
\input{sections/3_experiments_efficiency}
\input{figures/fig3_psnr_params}
\input{sections/3_experiments_visual}
\input{figures/fig7_visual}
\input{sections/3_experiments_ablation}
\input{tables/table_c1_dpr_vs_ema}
\input{tables/table3_ablation_a1}
\input{tables/table4_ablation_a2}
\input{tables/table5_ablation_a3}
\input{sections/3_experiments_routing_analysis}
% =====================================================================
% Fig.5 -- Per-token routing-confidence distribution x^s = max_m S_{n,m}.
% outline.md §3.6 (routing interpretability) / traceability C3 ("x_scores 分布").
% Self-contained pgfplots figure; \input into the paper.
% Requires in preamble: \usepackage{pgfplots}\pgfplotsset{compat=1.16}
%
% DATA PROVENANCE (verbatim, no fabrication):
%   fig58_assets/fig5_{Urban100,Manga109}_xscore_hist.csv produced on the
%   training machine by CATANet/scripts/build_routing_figures.py with the x4
%   model net_g_250000, accumulating x^s over ALL LR images of each test set
%   (Urban100: 4,839,460 tokens/block; Manga109: 6,545,996 tokens/block).
%   Curves below are the normalized frequency (% of tokens per 0.02 bin) for
%   three representative blocks b0 (M=16), b2 (M=64), b3 (M=128). Uniform
%   floor 1/M shown as dotted vertical lines (b0: 0.0625, b2: 0.0156,
%   b3: 0.0078). Per-block stats in fig5_*_xscore_stats.csv: block-mean x^s
%   sits at ~1.5-1.9x the floor (e.g. Urban100 b0 0.1093 vs 0.0625, b3 0.0124
%   vs 0.0078), i.e. informative but soft assignment. >99.8% of mass lies in
%   [0,0.30], so the x-axis is clipped there for readability.
% =====================================================================
\begin{figure}[H]
\centering
\pgfplotsset{
  xscoreaxis/.style={
    width=0.98\linewidth, height=4.2cm,
    xmin=0, xmax=0.30, ymin=0, ymax=100,
    ytick={0,25,50,75,100},
    xlabel={routing confidence $x^{\mathrm{s}}=\max_m S_{n,m}$},
    ylabel={tokens (\%)},
    ymajorgrids, grid style={gray!22},
    tick label style={font=\scriptsize},
    label style={font=\footnotesize},
    legend style={font=\scriptsize, at={(0.97,0.95)}, anchor=north east,
                  draw=none, fill=white, fill opacity=0.7, text opacity=1,
                  row sep=-1pt},
    title style={font=\footnotesize, yshift=-4pt},
  },
}
\begin{minipage}{\linewidth}
\centering
\begin{tikzpicture}
\begin{axis}[xscoreaxis, title={(a) Urban100}]
\addplot[draw=blue!70, very thick] coordinates {
  (0.07,16.2111) (0.09,31.6671) (0.11,23.1137) (0.13,15.4588) (0.15,6.6313)
  (0.17,2.9667) (0.19,1.5974) (0.21,0.9969) (0.23,0.397) (0.25,0.3825)
  (0.27,0.2386) (0.29,0.1443) };
\addplot[draw=orange!85!black, very thick] coordinates {
  (0.01,18.7685) (0.03,66.9281) (0.05,9.506) (0.07,2.6561) (0.09,1.0022)
  (0.11,0.4613) (0.13,0.2415) (0.15,0.1404) (0.17,0.089) (0.19,0.0578)
  (0.21,0.0375) (0.23,0.0282) (0.25,0.02) (0.27,0.015) (0.29,0.0116) };
\addplot[draw=teal, very thick] coordinates {
  (0.01,94.2813) (0.03,5.4999) (0.05,0.2043) (0.07,0.0134) (0.09,0.0009) };
\draw[blue!70, dotted, thick] (axis cs:0.0625,0) -- (axis cs:0.0625,100);
\draw[orange!85!black, dotted, thick] (axis cs:0.0156,0) -- (axis cs:0.0156,100);
\draw[teal, dotted, thick] (axis cs:0.0078,0) -- (axis cs:0.0078,100);
\legend{$b_0$ ($M{=}16$), $b_2$ ($M{=}64$), $b_3$ ($M{=}128$)}
\end{axis}
\end{tikzpicture}
\end{minipage}

\vspace{2pt}
\begin{minipage}{\linewidth}
\centering
\begin{tikzpicture}
\begin{axis}[xscoreaxis, title={(b) Manga109}]
\addplot[draw=blue!70, very thick] coordinates {
  (0.07,3.4849) (0.09,21.1291) (0.11,24.5368) (0.13,18.8673) (0.15,11.2919)
  (0.17,7.2655) (0.19,4.2136) (0.21,2.8329) (0.23,1.9641) (0.25,1.4542)
  (0.27,1.009) (0.29,0.6623) };
\addplot[draw=orange!85!black, very thick] coordinates {
  (0.01,11.53) (0.03,75.3017) (0.05,9.1121) (0.07,2.1833) (0.09,0.8387)
  (0.11,0.4047) (0.13,0.2143) (0.15,0.1256) (0.17,0.0812) (0.19,0.0557)
  (0.21,0.0398) (0.23,0.0267) (0.25,0.018) (0.27,0.0139) (0.29,0.0107) };
\addplot[draw=teal, very thick] coordinates {
  (0.01,95.4593) (0.03,4.1735) (0.05,0.3019) (0.07,0.051) (0.09,0.011)
  (0.11,0.0025) };
\draw[blue!70, dotted, thick] (axis cs:0.0625,0) -- (axis cs:0.0625,100);
\draw[orange!85!black, dotted, thick] (axis cs:0.0156,0) -- (axis cs:0.0156,100);
\draw[teal, dotted, thick] (axis cs:0.0078,0) -- (axis cs:0.0078,100);
\legend{$b_0$ ($M{=}16$), $b_2$ ($M{=}64$), $b_3$ ($M{=}128$)}
\end{axis}
\end{tikzpicture}
\end{minipage}
\caption{Distribution of the per-token routing confidence
$x^{\mathrm{s}}=\max_m S_{n,m}$ (Eq.~\eqref{eq:belong}) at $\times4$, accumulated
over \emph{all} test images of (a)~Urban100 and (b)~Manga109 with the final model.
Curves are shown for three representative TAB blocks with increasing prototype
count $M$; dotted vertical lines mark the uniform-assignment floor $1/M$. In every
plotted block the distribution sits clearly to the right of its floor (block means
are about $1.5$--$2.1$ times $1/M$) yet stays far from $1$, confirming that the softmax
assignment is \emph{informative but soft}: tokens commit to a prototype more than
chance while retaining the graded confidence that the C3 components
(ordering key, IASA gate, soft fallback) rely on. As $M$ grows the mode shifts
toward the (lower) floor, consistent with confidence being spread over more slots.}
\label{fig:xscore_hist}
\end{figure}

\input{figures/fig6_usage_entropy}
\input{figures/fig8_cluster_maps}
\input{sections/4_discussion}
\input{sections/5_conclusion}

%%%%%%%%% BACK MATTER (Springer style: declarations after conclusion)
\section*{Statements and Declarations}

\paragraph{Author contributions.}
Hainan Huang: conceptualization, methodology, software, validation, formal analysis,
investigation, data curation, writing -- original draft, visualization. Xiang Li:
conceptualization, supervision, funding acquisition, resources, writing -- review \&
editing, project administration. All authors have read and approved the submitted
version of the manuscript.

\paragraph{Funding.}
This research received no external funding.

\paragraph{Competing interests.}
The authors declare that they have no competing interests.

\paragraph{Data availability.}
The datasets used in this study are all publicly available: the DIV2K training set and the
Set5, Set14, BSD100, Urban100, and Manga109 benchmarks. No new data were generated. The code
and trained models are not publicly available due to institutional restrictions.

\paragraph{Use of generative AI.}
During manuscript preparation, generative AI tools were used to polish language, improve
readability, and assist in organizing the manuscript structure. The authors reviewed and
edited all resulting text and take full responsibility for the content of the publication.

%%%%%%%%% REFERENCES
\bibliography{references}

\end{document}
